\begin{frame}[t]{\ev{\tagO}Hedged requests reduce tail latency}
\vspace{0pt}
\lead{If a server is slow, ask a second server too, and use whichever answers first.}
\vspace{.05cm}
\begin{center}
{\fontsize{40}{46}\selectfont\bfseries 1,800 ms $\to$ 74 ms\par}
\vspace{.15cm}
{\large Latency of the slowest 0.1\% of requests. The backup requests add 2\% to the load.\par}
\vspace{.25cm}
{\small\color{Muted}``the source of latency is often not inherent in the particular request'': the server is slow, not the request.\par}
\end{center}
\vspace{.1cm}
\pause
\claim{Agents hedge for correctness. A flaw in the task repeats in every copy.}
\sources{Dean \& Barroso, The Tail at Scale, CACM 2013, backup after 10 ms, 99.9th-percentile latency reading 1,000 BigTable keys. It works only when the slowdown does not hit several servers at once. Shared mistakes across different AI models: \href{https://arxiv.org/abs/2506.07962}{Kim et al., ICML 2025}.}
\note{[1:10] In 2013 Jeff Dean and Luiz Barroso published a technique called hedging. The client sends a request to one server, and if no answer arrives within 10 milliseconds it sends a copy to a second server and uses the first answer. In one benchmark, the slowest requests fell from 1,800 milliseconds to 74, for two percent more requests. The delay usually originates in the server rather than in the request. Two servers rarely stall together. [click] Agents now hedge for correctness by running the task several times and keeping a run that passes. When the flaw is in the task itself, every copy repeats it. Their paper describes a second technique for that case, presented at the end.}
\end{frame}
